\subsection{Priors, Posteriors, and Conjugacy}

Bayesian inference treats unknown parameters as uncertain quantities and updates a prior distribution by the likelihood supplied by the data.

\begin{definition} \label{def:prior-posterior-predictive} \hlt{Prior, Posterior, and Predictive Distribution}\\
Let $\Pi$ be a prior distribution on $\Theta$ and let $p_\theta(x)$ be a dominated likelihood.  The \hlt{posterior} is
\[
d\Pi(\theta\mid x)
=\frac{p_\theta(x)\,d\Pi(\theta)}
{\int_\Theta p_\vartheta(x)\,d\Pi(\vartheta)}
\]
when the denominator is finite and positive.  The \hlt{posterior predictive distribution} for a future observation $X_{\mathrm{new}}$ is
\[
\mathbb{P}(X_{\mathrm{new}}\in A\mid x)
=\int_\Theta P_\theta(A)\,d\Pi(\theta\mid x).
\]
\end{definition}

\begin{definition} \label{def:conjugate-prior-family} \hlt{Conjugate Prior}\\
A prior family $\{\Pi_\eta:\eta\in H\}$ is \hlt{conjugate} for a likelihood family if the posterior belongs to the same family after observing data.  Conjugacy turns Bayes updating into hyperparameter updating.
\end{definition}

\begin{example} \label{ex:beta-binomial-conjugacy} \hlt{Beta--Binomial Updating}\\
If $X\mid p\sim\operatorname{Binomial}(n,p)$ and $p\sim\operatorname{Beta}(a,b)$, then
\[
p\mid X=x\sim\operatorname{Beta}(a+x,b+n-x).
\]
The posterior mean is $(a+x)/(a+b+n)$, a weighted average of the prior mean $a/(a+b)$ and the sample proportion $x/n$.
\end{example}

\begin{example} \label{ex:normal-normal-conjugacy} \hlt{Normal Mean with Normal Prior}\\
Let $X_1,\ldots,X_n\mid\mu$ be independent $\mathcal{N}(\mu,\sigma^2)$ with known $\sigma^2$, and let $\mu\sim\mathcal{N}(m_0,\tau_0^2)$.  Then
\[
\mu\mid X
\sim\mathcal{N}(m_n,\tau_n^2),
\qquad
\tau_n^{-2}=\tau_0^{-2}+n\sigma^{-2},
\]
and
\[
m_n=\tau_n^2\left(\tau_0^{-2}m_0+n\sigma^{-2}\overline{X}\right).
\]
\end{example}

\begin{remark} \label{rem:bayesian-conditioning-measure-theory} \hlt{Conditioning Viewpoint}\\
Bayesian updating is ordinary conditioning on a joint probability space for $(\theta,X)$.  The measure-theoretic conditional expectation and disintegration ideas from earlier probability parts supply the formal foundation.
\end{remark}
